% ======================================================================
\section{Open challenges and outlook}
\label{sec:outlook}
% ======================================================================

\paragraph{Scaling to high-Reynolds-number, three-dimensional flows.}
Most demonstrations of learned representations use two-dimensional or
low-Reynolds-number flows, for which global bases or small autoencoders
suffice. Operationally relevant flows are different in kind. Ship airwakes,
for example, are massively separated flows at Reynolds numbers of order
$10^{7}$--$10^{8}$ that must be predicted over many wind-over-deck conditions
for flight simulation and flight-deck safety analysis
\parencite{polsky2016safedi}. At this scale even storing global snapshots
strains memory, dense encoders are infeasible, and the number of
distinct flow patterns grows with domain size and Reynolds number.
Representations whose cost and training data grow with the \emph{local}
complexity of the flow rather than its global complexity are therefore
attractive. These include domain-decomposed and GFEM-enriched bases
\parencite{hoang2021ddlspg,diaz2024ddnm,deshmukh2020dls}, graph and
local-token architectures \parencite{barwey2023multiscale,serrano2024aroma},
and in situ compression \parencite{glaws2020insitu}. Whether forecasting
itself can be posed locally without losing long-range coupling is an open
question. So is how locality interacts with the global pressure field of
incompressible flow.

\paragraph{Generalization across parameters and conditions.}
Interpolating between training parameters works well for smooth parametric
dependence \parencite{amsallem2008grassmann,fries2022lasdi}, but extrapolation to
new geometries, Reynolds numbers, or operating conditions remains the weakest
point of learned representations. Local bases may generalize better because
local flow features (shear layers, boundary layers, vortices) recur across
conditions even when their global arrangement changes. This hypothesis
deserves systematic testing.

\paragraph{Experimental and sparse data.}
Experimental data are noisy, gappy, and often not time-resolved. Classical
tools include gappy POD \parencite{everson1995gappy}, spectral estimation from
non-time-resolved PIV synchronized with pressure sensors
\parencite{zhang2020samm}, and sparse-sensor reconstruction
\parencite{manohar2018sensor,callaham2019robust,erichson2020shallow}. They need to
be integrated with learned representations, for example by training encoders
that accept partial observations or by using generative latent models for
conditional reconstruction \parencite{du2024cnf}.

\paragraph{Uncertainty, statistics, and generative modeling.}
For chaotic flows, point forecasts beyond a few Lyapunov times are
meaningless. What matters is whether a model reproduces the statistics and the
distribution of future states. Probabilistic latent spaces and latent
diffusion models \parencite{kohl2023acdm,gao2024bayesian,du2024cnf} are promising
but expensive, and their evaluation calls for statistical and spectral metrics
rather than pointwise error.

\paragraph{Interpretability and physical consistency.}
Linear and local bases offer interpretable modes and simple error estimates.
Learned representations offer compactness at the cost of both. Continued work
on ordered and disentangled latent variables
\parencite{fukami2020hierarchical,eivazi2022betavae}, on relating latent
coordinates to invariant solutions and physical mechanisms
\parencite{page2021revealing,farzamnik2023snapshots}, and on hard physical
constraints \parencite{lee2021deepconservation,mohan2020embedding} will be needed
before learned representations are trusted in engineering practice.

\paragraph{Foundation models and tokenization.}
Pretrained models for PDEs \parencite{mccabe2024mpp,herde2024poseidon} and large
multi-physics data collections \parencite{ohana2024well} raise a new version of
the representation question: how should heterogeneous flow data be tokenized?
Patch tokens \parencite{dosovitskiy2021vit}, learned slices
\parencite{wu2024transolver}, latent tokens \parencite{alkin2024upt}, and
physically structured local bases \parencite{deshmukh2020dls} embody different
answers. The DR literature reviewed here, with its emphasis on optimality,
locality, and multi-scale structure, is directly relevant to that choice.

% ======================================================================
\section{Conclusions}
\label{sec:conclusions}
% ======================================================================

Dimensionality reduction is the foundation on which data-driven models of
fluid flows are built. It fixes the information available to the dynamics
model, the cost of training and inference, and the difficulty of the learning
problem. Linear modal decompositions remain the indispensable baseline: they are
optimal among linear subspaces, cheap, interpretable, and, for flows dominated
by a few coherent structures, hard to beat. Their limitations for transport and
multi-scale turbulence are well understood through the Kolmogorov $n$-width.
Two broad responses have emerged. Nonlinear manifold methods, from kernel and
graph embeddings to deep and variational autoencoders and implicit neural
representations, achieve compact latent spaces when the intrinsic dimension is
low, at the cost of data, tuning, and an accuracy ceiling. Localized
representations, from domain decomposition to GFEM enrichment and data-informed
local subspaces, keep linear, sparse encoders and decoders with guaranteed
continuity. They learn from few snapshots and scale with the local rather than
global complexity of the flow, at the cost of larger latent spaces when the
flow is globally simple.

The literature increasingly suggests that these choices are complementary,
and that representations should be judged by how well they support learning
the dynamics: whether the latent state is sufficient, smooth, stable, and
statistically faithful, not only by reconstruction error. Progress toward
high-Reynolds-number, three-dimensional, and parametrically varying flows of
engineering interest will depend on representations that respect the
multi-scale and local structure of turbulence. It will also depend on
evaluation against strong baselines, held-out conditions, and statistical
metrics.
